\begin{definition}[Турнир]
Турнир — орграф, в котором между любыми двумя различными вершинами
есть ровно одна дуга одного из двух направлений; петель нет.
Гамильтоновы путь и цикл проходят через все вершины по направлениям дуг.
\end{definition}
\begin{lemma}[Вставка в путь]
Если \(v_1\to\cdots\to v_k\) — путь в турнире и \(x\) — новая вершина,
то её можно вставить в этот путь с сохранением порядка старых вершин.
\end{lemma}
\begin{proof}
Если \(x\to v_1\), добавим \(x\) в начало; если \(v_k\to x\), в конец.
Иначе \(v_1\to x\) и \(x\to v_k\). Пусть \(j\) — первый индекс,
для которого \(x\to v_j\). Тогда \(j\ge2\), а \(v_{j-1}\to x\),
поэтому вставка между \(v_{j-1}\) и \(v_j\) корректна.
\end{proof}
\begin{theorem}[Редеи]
В любом непустом турнире существует гамильтонов путь.
\end{theorem}
\begin{proof}
Начинаем с одной вершины и вставляем остальные по лемме.
После каждой вставки путь содержит все уже обработанные вершины;
через конечное число шагов он содержит весь турнир.
\end{proof}
\begin{lemma}[Вставка в цикл]
Если вершина \(x\) вне направленного цикла имеет как входящие,
так и исходящие дуги к его вершинам, её можно вставить в цикл.
\end{lemma}
\begin{proof}
При циклическом просмотре вершин есть переход от типа
\(v_i\to x\) к типу \(x\to v_{i+1}\), поскольку оба типа присутствуют.
Заменим дугу \(v_i\to v_{i+1}\) двумя дугами через \(x\).
\end{proof}
\begin{graphillustration}{Вставка $x$ заменяет одну дугу двумя. На схеме цикла показан пример с тремя исходными вершинами; на пути возможна также вставка на конце.}{fig:12-insert}
\begin{tikzpicture}[font=\small]
\begin{scope}[shift={(0,0)}]
\node[gi vertex] (p) at (0,0) {$v_{j-1}$};
\node[gi vertex] (q) at (3.5,0) {$v_j$};
\node[gi selected] (x) at (1.75,1.3) {$x$};
\draw[gi removed,-{Stealth}] (p) -- (q);
\draw[gi arrow] (p) -- (x);
\draw[gi arrow] (x) -- (q);
\node[gi note] at (-0.7,0) {$\cdots$};
\node[gi note] at (4.2,0) {$\cdots$};
\node[gi note] at (1.75,-1.9) {Вставка внутри пути};
\end{scope}
\begin{scope}[shift={(6.5,0)}]
\node[gi vertex] (p) at (0,0) {$v_i$};
\node[gi vertex] (q) at (3.5,0) {$v_{i+1}$};
\node[gi selected] (x) at (1.75,1.3) {$x$};
\draw[gi removed,-{Stealth}] (p) -- (q);
\draw[gi arrow] (p) -- (x);
\draw[gi arrow] (x) -- (q);
\node[gi vertex] (w) at (1.75,-1.1) {$w$};
\draw[gi arrow] (q) -- (w);
\draw[gi arrow] (w) -- (p);
\node[gi note] at (1.75,-1.9) {Вставка в цикл};
\end{scope}
\end{tikzpicture}
\end{graphillustration}


\begin{theorem}[Камион; критерий Редеи–Камиона]
Турнир с \(n\ge3\) имеет гамильтонов цикл тогда и только тогда,
когда он сильно связен.
\end{theorem}
\begin{proof}
Гамильтонов цикл даёт направленный путь между любой упорядоченной парой,
значит, сильная связность необходима.
Пусть турнир сильно связен. В нём есть направленный цикл:
возьмём дугу \(u\to v\) и простой обратный путь из \(v\) в \(u\);
его длина не меньше двух, поскольку противоположной дуги нет.
Выберем простой цикл \(C\) максимальной длины.
Если \(C\) не гамильтонов, по лемме каждая внешняя вершина либо
получает дуги от всех вершин \(C\), либо посылает дуги ко всем ним.
Обозначим эти множества \(A\) и \(B\) соответственно.

Оба множества непусты: при \(A=\varnothing\) из цикла нельзя выйти,
а при \(B=\varnothing\) нельзя в него попасть, вопреки сильной связности.
Есть дуга \(a\to b\), \(a\in A,b\in B\): иначе все дуги между ними
направлены из \(B\) в \(A\) и из \(A\) нельзя достичь цикла.
Теперь любую дугу \(v_i\to v_{i+1}\) цикла можно заменить путём
\(v_i\to a\to b\to v_{i+1}\), увеличив цикл на две вершины.
Противоречие максимальности. Поэтому \(C\) гамильтонов.
\end{proof}

\begin{graphillustration}{Шаг доказательства Камиона: внешним путём $v_i\to a\to b\to v_{i+1}$ заменяем дугу цикла. Схема показывает только используемые дуги турнира.}{fig:12-camion}
\begin{tikzpicture}[font=\small]
\node[gi vertex] (vi) at (0,0) {$v_i$};
\node[gi vertex] (vj) at (4.5,0) {$v_{i+1}$};
\node[gi vertex] (w) at (2.25,-1.3) {$w$};
\node[gi vertex] (a) at (1,1.5) {$a\in A$};
\node[gi vertex] (b) at (3.5,1.5) {$b\in B$};
\draw[gi arrow] (vi) -- (a);
\draw[gi arrow] (a) -- (b);
\draw[gi arrow] (b) -- (vj);
\draw[gi arrow] (vj) -- (w);
\draw[gi arrow] (w) -- (vi);
\draw[gi removed,-{Stealth}] (vi) -- (vj);
\node[gi note] at (7.7,0.5) {$C\to A,\quad A\ni a\to b\in B,\quad B\to C$\\[5pt]Цикл удлиняется на две вершины.};
\end{tikzpicture}
\end{graphillustration}

